% §2 What the line paper settled --- the recap of the released line paper
% (\sscite{fagerstrom2026spatial}, v0.1, concept DOI 10.5281/zenodo.22761631): the notation
% table, the axioms, the exponent and its profile, the characterization theorem, and the one
% chapter-2 node of the blueprint that paper did not print.
%
% Rewritten 2026-09-16 at the author's review. The first draft (2026-09-15) cited everything and
% restated nothing, "the minimum the later sections need to resolve"; the author's decision is
% that this paper is readable on its own, so what the later sections read is restated here:
% the statements verbatim from the blueprint (the text of record, `% shared with blueprint`
% markers, guarded by `linkage check`), each with the line paper's statement number in its
% title and its proof left in the line paper; the rest --- the gauge, the folded profile, the
% symbol, the hemigroup --- in prose. Restated labels are local and are commented out of the
% line-number table in macros.tex; every other \ref to the line paper still prints "n of [line]".
% The notation table is the programme's shared header (hub RELEASES.md rule 3), reproduced with
% the bridge rows added, which is a versioned change recorded in the hub log; the line paper
% picks the rows up at its next release. Nothing here is proved except the last proposition.
%
% Contents:
%   2.1 notation (spaces, actions, convolution; Table 2);
%   2.2 the axioms (def:cascade-family; a stage; hemigroup versus semigroup);
%   2.3 exponents and profiles (the Levy-Khintchine form, self-decomposability, prop:sd-exponents,
%       the folded profile, lem:profile-integrability, admissible pairs, the displacement
%       dictionary copied from the line paper's Remark 2.6, the symbol B);
%   2.4 the characterization theorem (the canonical gauge; thm:main-characterization,
%       cor:semigroup-case, lem:admissible-cone, prop:kernel-regularity);
%   2.5 what the line paper takes on trust;
%   2.6 prop:laplace-uniqueness-locally-finite (blueprint chapter 2), [T] and \leanok on Lean
%       core, with its proof: the line paper left it out because its consumers are here.

\section{What the line paper settled}\label{sec:line}

This paper continues the line paper \sscite{fagerstrom2026spatial}, and is written to be read
on its own. The main things the later sections use of that paper are restated in this section:
its axioms, the form of its exponents, and its characterization theorem, together with the
notation. A restated statement carries the line paper's statement number in its title, and its
proof is at that number. Further results are used and only cited, in the form
``Lemma~\ref{lem:selfdecomposable-exponents}'' with the number of that paper's first release.
Four are used throughout:
the Fourier toolbox (Bochner's theorem and the \Levy--Khintchine representation,
Proposition~\ref{prop:fourier-toolbox}), the
quadratic growth of an exponent (Lemma~\ref{lem:quadratic-growth}), the uniqueness of the
transform of a measure (Proposition~\ref{prop:fourier-uniqueness}), and \Levy's
continuity theorem (Proposition~\ref{prop:levy-continuity}). The others are used once or twice
each: the differentiability of a self-decomposable exponent and its representation
(Lemma~\ref{lem:selfdecomposable-exponents}, in \S\S\ref{sec:generator}
and~\ref{sec:implementation}), the two members that are not stable
(Proposition~\ref{prop:two-members}, in \S\S\ref{sec:bridge} and~\ref{sec:corners}), the zero
set of an exponent and the absence of lattice kernels (Lemmas~\ref{lem:lattice-zero}
and~\ref{lem:no-lattice}, in \S\ref{sec:bridge}), and the additivity of the stage exponents with
the nonvanishing of the transforms (Lemmas~\ref{lem:additivity} and~\ref{lem:nonvanishing}, in
\S\ref{sec:implementation}). One classical inequality is used from inside a proof of the line
paper, that of Theorem~\ref{thm:increments-levy}, and is restated here because
Lemma~\ref{lem:cin-delay-equation} rests on it. For a symmetric probability measure $\mu$ and
$\eta > 0$,
\begin{equation}\label{eq:truncation-b}
  \int \Bigl(1 - \frac{\sin\eta x}{\eta x}\Bigr)\,\mu(dx)
   = \frac1\eta\int_0^\eta \bigl(1 - \hat\mu(\omega)\bigr)\,d\omega,
  \qquad
  1 - \frac{\sin u}{u} \ge \frac{2}{3\pi^2}\,\bigl(1 \wedge u^2\bigr),
\end{equation}
the first by Fubini's theorem and the second elementary. Together they give the truncation
inequality $\mu\{|x| \ge 1/\eta\} \le \tfrac{3\pi^2}{2}\,\eta^{-1}\int_0^\eta(1 - \hat\mu)$. When
$\hat\mu = e^{-F}$ with $F(\omega) \le C\omega^2$ near the origin, the right-hand side is
$O(\eta^2)$, since $1 - e^{-F} \le F$, so $\mu\{|x| > R\} = O(R^{-2})$ and $\mu$ has a finite
first absolute moment. The section ends with
one elementary proposition of the line paper's preliminaries that the line paper did not
print, because its uses are in this module; that one is proved.

\subsection{Notation}\label{sec:notation}

Space runs over $\RR$ with coordinate $x$, the wavenumber is $\omega$, and scale is $t$.
Table~\ref{tab:notation} is the notation table shared by the spatial modules, the line paper's
Table~2 with the
letters of the bridge added; the letters of \S\S\ref{sec:corners}--\ref{sec:implementation}
are introduced where they are used.

\emph{Spaces.} $\Lone = L^1(\RR, dx)$, real, with positive cone $L^1_+$ and
$\int f := \int_\RR f(x)\,dx$; $M_+(\RR)$ is the finite positive Borel measures on $\RR$, of
total mass $\|\mu\| = \mu(\RR)$. \emph{Actions.} Translations are $(\Tr{a} f)(x) = f(x-a)$.
The reflection is $(Rf)(x) = f(-x)$, extended to measures as the pushforward under
$x \mapsto -x$, and a measure is \emph{symmetric} when $R\mu = \mu$. Dilations act on
functions by $(\Dil{\lambda} f)(x) = \lambda^{-1} f(x/\lambda)$, $\lambda > 0$, which
preserves mass, and on measures by the pushforward under $x \mapsto \lambda x$.
\emph{Convolution and transform.} $(\mu * f)(x) = \int f(x-y)\,\mu(dy)$ for $\mu \in M_+(\RR)$
and $f \in \Lone$. Both actions are compatible with convolution. The transform of a symmetric
measure is $\hat\mu(\omega) = \int \cos(\omega x)\,\mu(dx)$, real and even, and the transform
of a stage's kernel is written $\hat\mu_{s,t}$. It is the restriction to symmetric measures of
$\hat\rho(\omega) = \int e^{-i\omega x}\,\rho(dx)$, the convention for a finite measure, positive
or signed, that is not symmetric. The sign matters in one place, the proof of
Lemma~\ref{lem:cin-delay-equation}, where a translate has the transform
$e^{i\omega\tau}\hat\nu(\omega)$, and the uniqueness of the transform
(Proposition~\ref{prop:fourier-uniqueness}) is stated in the line paper for every finite Borel
measure on $\RR$ in this convention.

\begin{table}[!htb]
\centering
\caption{Notation: the symbol, where it is defined, and the letter the causal article
uses for the same object. The rows above the rule are the line paper's Table~2; the rows
below it are the letters of the bridge, \S\ref{sec:bridge}, where the causal article's
Laplace variable is written $\sigma$, $s$ being a scale here. Brownian motion is written $\BM$, the letter $B$ being
the symbol of the exponent.}
\label{tab:notation}
\footnotesize
\setlength{\tabcolsep}{3pt}
\begin{tabular}{@{}lll@{}}
\toprule
object & here & the causal article \\
\midrule
signal coordinate & $x$ (space) & $t$ (time) \\
transform variable & $\omega$ (Fourier) & $s$ (Laplace) \\
scales in a stage & $r \le s \le t$ & $x \le y \le z$ \\
stage operator and its kernel & $\Phis{s}{t}$, $\mu_{s,t}$ & $\Phi_{x,y}$, $\mu_{x,y}$ \\
exponent of a stage & $g_{s,t}(\omega) = -\log\hat\mu_{s,t}(\omega)$ & $g_{x,y}(s)$ \\
kernel from the origin, accumulated exponent & $\mu_{0,t}$, $G(t,\omega) = g_{0,t}(\omega)$ & $\mu_{0,x}$, $G(x,s)$ \\
canonical gauge & $\chi(t)$ (\S\ref{sec:characterization-recap}) & $\chi(x)$ \\
the family's exponent and its pair & $F(\omega)$; Gaussian coefficient $a$, profile $k$ & $F(s)$; drift $b_0$, profile $k$ \\
\Levy{} measure, tail measure of the profile & $\nu$ \ref{eq:levy-khintchine}, $\varpi = -dk$ \ref{eq:symbol} & $\nu$, $\varpi$ \\
symbol of the exponent & $B(\omega) = \omega F'(\omega)$ \ref{eq:symbol} & $B(s) = sF'(s)$ \\
dilation ratio, action, orbit coordinate & $\lambda$, $S_\lambda$, $c(\lambda)$ & $\sigma$, $S_\sigma$, $c(\sigma)$ \\
random displacement, random delay & $X_{s,t}$ & $T_{x,y}$ \\
function space & $\Lone = L^1(\RR,dx)$ & $X = L^1(\RR,dt)$ \\
translation, reflection & $\Tr{a}$, $R$ & $\tau_a$, --- \\
exponents by definition, by representation & $\NDs$, $\LEs$ (\S\ref{sec:exponents-recap}) & $\mathrm{BF}_0$, $\mathrm{LE}$ \\
\midrule
standard Brownian motion, its law at time $u$ & $\BM$, $g_u$ & --- \\
the causal exponent and its data & $F_{\mathrm I}(\sigma)$; $(b_0, k_{\mathrm I})$ (Def.~\ref{def:causal-admissible}) & $F(s)$; $(b_0, k)$ \\
the random delay of the bridge & $T_1$, with exponent $F_{\mathrm I}$ & $T_{0,1}$ \\
\bottomrule
\end{tabular}
\end{table}

\subsection{The axioms}\label{sec:axioms-recap}

The line paper's object is a family of operators indexed by pairs of scales. A pair of scales
$s \le t$ is a \emph{stage}: the operator $\Phis{s}{t}$ of the stage takes the signal at scale
$s$ to the signal at scale $t$, and its kernel $\mu_{s,t}$, which the line paper's
representation lemma supplies (its Lemma~\linenum{lem:convolution-representation}), is the
law of the random displacement the stage adds. The axioms are these.

% shared with blueprint def:cascade-family
\begin{definition}[symmetric cascade measurement family, Definition~3.1 of the line paper]\label{def:cascade-family}
A family of linear operators $\Phis{s}{t} : \Lone \to \Lone$, $0 \le s \le t < \infty$, is a
\emph{symmetric cascade measurement family} if:
\begin{enumerate}
\item[(A1)] \emph{Boundedness.} Each $\Phis{s}{t}$ is a bounded operator on $\Lone$.
\item[(A2)] \emph{Translation covariance.} $\Phis{s}{t}\,\Tr{a} = \Tr{a}\,\Phis{s}{t}$ for
every $a \in \RR$.
\item[(A3)] \emph{Reflection symmetry.} $\Phis{s}{t}\,R = R\,\Phis{s}{t}$.
\item[(A4)] \emph{Positivity.} $\Phis{s}{t}\,L^1_+ \subseteq L^1_+$.
\item[(A5)] \emph{Unit mass.} $\int \Phis{s}{t} f = \int f$ for every $f \in L^1_+$.
\item[(A6)] \emph{Cascade (hemigroup).} $\Phis{t}{t} = \mathrm{Id}$, and
$\Phis{s}{t}\,\Phis{r}{s} = \Phis{r}{t}$ for all $r \le s \le t$.
\item[(A7)] \emph{Continuity.} For every $f \in \Lone$ the map $(s,t) \mapsto \Phis{s}{t}f$ is
continuous from $\{0 \le s \le t\}$ to $\Lone$.
\item[(A8)] \emph{Scale covariance.} There is a family $(S_\lambda)_{\lambda > 0}$ of
increasing bijections of $[0,\infty)$ with
$\Dil{\lambda}\,\Phis{s}{t} = \Phi_{S_\lambda s,\,S_\lambda t}\,\Dil{\lambda}$ for all
$0 \le s \le t$ and all $\lambda > 0$.
\item[(ND)] \emph{Nondegeneracy.} $\Phis{s}{t} \ne \mathrm{Id}$ whenever $s < t$.
\end{enumerate}
\end{definition}

A family satisfying the nine is called \emph{admissible} throughout this paper. The word is
also used of a pair $(a,k)$ and of an exponent, in the senses fixed in
\S\ref{sec:exponents-recap} and after Theorem~\ref{thm:main-characterization}.

\emph{Hemigroup versus semigroup.} The axiom that separates this class from the classical one
is (A6). The scale-space axioms of \sscite{pauwels1995extended} ask for a one-parameter family
with $K_s K_t = K_{s+t}$, so that a stage depends on the difference $t - s$ of its scales
only: the family is a convolution semigroup, and the kernel at scale $t$ is the law at time
$t$ of a process with stationary independent increments. Axiom (A6) asks only that the stages
compose. The stage from $r$ to $t$ is the stage from $s$ to $t$ after the stage from $r$
to $s$; how the stage from $s$ to $t$ depends on the two scales is left free. The family is
then the family of increment laws of a process with independent increments that need not be
stationary in scale, and scale covariance (A8) fixes what those increments are:
the process is self-similar, and its marginals are self-decomposable. That is the difference
in the composition axiom, and it is the reason the class here has a profile $k$
where the semigroup class has one index $\alpha$
(Corollary~\ref{cor:semigroup-case}). The hemigroup form of the axioms comes from the
time-causal article \sscite{fagerstrom2026hemigroup}, whose stages run over time on the
half-line; the line paper carried them to the reflection-symmetric line, and
\S\ref{sec:bridge} is where the two meet.

\subsection{Exponents and profiles}\label{sec:exponents-recap}

Under the axioms, every stage is a convolution, $\Phis{s}{t} f = \mu_{s,t} * f$, with a
symmetric probability measure $\mu_{s,t}$ whose transform does not vanish, and the
\emph{exponent} of the stage is $g_{s,t} = -\log\hat\mu_{s,t}$ (the line paper's
Lemmas~\linenum{lem:convolution-representation} and~\linenum{lem:nonvanishing}). The
exponents add along the cascade, and each is a symmetric continuous negative definite
function. Such a function is a continuous even $\psi \ge 0$ with $\psi(0) = 0$ for which
$e^{-\tau\psi}$ is the transform of a symmetric probability measure for every $\tau > 0$.
They form the class $\NDs$ of the line paper's Definition~\linenum{def:symmetric-negdef},
whose members are exactly the functions below. This is the symmetric \Levy--Khintchine
representation, clause~(3) of Proposition~\ref{prop:fourier-toolbox}, which asserts the
representation, its converse and the uniqueness of the pair:
\begin{equation}\label{eq:levy-khintchine}
  \psi(\omega) \;=\; a\,\omega^2 + \int_{(0,\infty)} \bigl(1 - \cos\omega x\bigr)\,\nu(dx),
  \qquad a \ge 0, \quad \int_{(0,\infty)} (1 \wedge x^2)\,\nu(dx) < \infty,
\end{equation}
with the pair $(a,\nu)$ unique. The line paper names the class by this representation,
$\LEs$ (its Definition~\linenum{def:levy-exponents}), so that $\NDs = \LEs$, and this paper
uses both names for the one class. In the representation $a$ is the
\emph{Gaussian coefficient} and $\nu$ the \emph{\Levy{} measure}. The measure $\nu$ lives on
the positive half-line by a convention that matters for constants: the \Levy{} measure of a
symmetric law on $\RR$ is a symmetric measure $\nu_2$ on $\RR \setminus \{0\}$, and $\nu$ is
its image under $x \mapsto |x|$, the two directions of a displacement of a given size counted
together.

A symmetric law $\mu$ is \emph{self-decomposable} (the line paper's
Definition~\linenum{def:self-decomposable}) if for every $0 < c < 1$ there is a law $\rho_c$ with
$\mu = (\Dil{c}\mu) * \rho_c$: a random displacement with law $\mu$ is, for every contraction
ratio, a contracted copy of itself plus an independent residual. The characterization theorem
below says that the kernels of an admissible family are of this kind, and this is the class
that has a profile.

% shared with blueprint prop:sd-exponents
\begin{proposition}[symmetric self-decomposable exponents, Proposition~2.13 of the line paper]\label{prop:sd-exponents}
A symmetric infinitely divisible law with pair $(a,\nu)$ is self-decomposable if and only if
its \Levy{} measure has the form $\nu(dx) = k(x)\,x^{-1}dx$ on $(0,\infty)$ with
$k : (0,\infty) \to [0,\infty)$ nonincreasing, so that its exponent is
\begin{equation}\label{eq:sd-profile}
  \psi(\omega) \;=\; a\,\omega^2
  + \int_0^\infty \bigl(1 - \cos\omega x\bigr)\,\frac{k(x)}{x}\,dx,
  \qquad a \ge 0, \quad k \ge 0 \ \text{nonincreasing}.
\end{equation}
Self-decomposability imposes no restriction on the Gaussian coefficient $a$.
\end{proposition}

The function $k$ is the \emph{folded profile} of the law, or its profile for short. It is
folded because $\nu$ is: a two-sided \Levy{} density written $h(|x|)/|x|$, the form of
\sscite{sato1999levy}, Cor.~15.11, p.~95, has folded profile $k = 2h$, and the factor is the
one place a
constant is easily lost in carrying a statement from a source into this paper. The profile
is the object of this paper. It is a nonincreasing function on the half-line, and the
representation reads it as an intensity: displacements of size about $x$ arrive at rate
$k(x)\,x^{-1}dx$, so $k$ is that rate per unit of $\log x$, and self-decomposability says
that the intensity of a displacement never grows with its size. The Gaussian is the member
with no profile, $k = 0$; every other member is described by the shape of $k$, and
\S\ref{sec:cone} describes the set of all shapes, \S\ref{sec:bridge} the shapes that come
from a causal delay, \S\ref{sec:corners} the shapes of the named families.

The integrability condition in \ref{eq:levy-khintchine}, read for a profile, is the following
criterion.

% shared with blueprint lem:profile-integrability
\begin{lemma}[the integrability criterion for a profile, Lemma~2.15 of the line paper]\label{lem:profile-integrability}
Let $k : (0,\infty) \to [0,\infty)$ be measurable and put $\nu(dx) := k(x)\,x^{-1}dx$. Then
\[
  \int_{(0,\infty)} (1 \wedge x^2)\,\nu(dx) < \infty
  \qquad\Longleftrightarrow\qquad
  \int_0^1 x\,k(x)\,dx < \infty \ \ \text{and} \ \
  \int_1^\infty \frac{k(x)}{x}\,dx < \infty .
\]
Consequently, by Proposition~\ref{prop:fourier-toolbox}(3), the right-hand pair of conditions
is necessary and sufficient for \ref{eq:sd-profile} to define a member of $\NDs$.
\end{lemma}

A pair $(a,k)$ with $a \ge 0$, $k : (0,\infty) \to [0,\infty)$ nonincreasing and the two
integrals of the criterion finite is called \emph{admissible}, and so is the exponent
\ref{eq:sd-profile} it defines; profiles equal almost everywhere define the same exponent.
The zero exponent is admissible as an exponent, and it is the one exponent that (ND) excludes
as a family.

The representation reads a kernel as a Gaussian part plus a catalogue of jumps. The reading is
the line paper's Remark~2.6, reproduced here because the cone section leans on it.

% line paper rem:displacement-dictionary (Remark 2.6 of v0.1, the number written out above
% since the label is local here), copied; the blueprint's remark
% of that label is the source of both halves of the line paper's split, and the paper's first
% half is what is reproduced here.
\begin{remark}[the displacement dictionary]\label{rem:displacement-dictionary}
A symmetric kernel is the law of a \emph{random displacement}: $\Phis{s}{t}f$ reports the
signal at a randomly displaced position $x - X_{s,t}$ with $X_{s,t} \sim \mu_{s,t}$. Its
transform $\hat\mu_{s,t}(\omega) = \Ex\cos(\omega X_{s,t})$ is the contrast that a
sinusoidal probe of wavenumber $\omega$ retains after the random displacement, the stage's
modulation transfer function (Remark~\ref{rem:modulation-transfer}). The exponent
$\psi = -\log\hat\mu$ adds along the cascade. In \ref{eq:levy-khintchine} the Gaussian
coefficient $a$ is the \emph{diffusive part}, a Gaussian displacement of variance $2a$, the
one term the time-causal axis deletes and
the only one Gaussian scale space uses. The measure $\nu$ is a \emph{catalogue of jump
displacements}. Three objects describe it in this paper, and they are kept apart by name: the
catalogue itself, the measure $\nu(dx) = k(x)\,dx/x$ on the half-line; the \emph{profile} $k$,
its density against $dx/x$; and the \emph{tail measure} $\varpi = -dk$ of \S\ref{sec:cone},
which says how much of each unit step the profile contains. Displacements of size about $x$
arrive as a Poisson stream, that is at
independent random moments with a constant rate, and $\nu(dx)$ is that rate; the streams of
different sizes are independent, and each event of size $x$ costs the probe the phase
deficit $1 - \cos\omega x$. So a symmetric infinitely divisible displacement is a Gaussian
part plus an independent sum of jumps of all sizes, and the representation says that every
such displacement decomposes in this way.
\end{remark}

\emph{The symbol.} An admissible exponent, written $F$ from here on, is $C^1$ away from the
origin, and the line paper
writes its derivative through the \emph{symbol}
\begin{equation}\label{eq:symbol}
  B(\omega) = \omega F'(\omega)
   = 2a\,\omega^2 + \int_{(0,\infty)} \bigl(1 - \cos\omega u\bigr)\,\varpi(du),
\end{equation}
where $\varpi = -dk$ is the \emph{tail measure} of the profile, the Stieltjes measure of the
nonincreasing function $k$ (the line paper's Lemma~\linenum{lem:selfdecomposable-exponents}).
The symbol is the
rate at which the exponent grows along a dilation, $B(\omega) = \partial_\lambda
F(\lambda\omega)|_{\lambda = 1}$, and it is itself an exponent of the form
\ref{eq:levy-khintchine}, with pair $(2a,\varpi)$. It is the object the scale evolution is
written in. In the canonical gauge of \S\ref{sec:characterization-recap} the kernel at scale
$t$ has transform $e^{-F(t\omega)}$, whose logarithmic derivative in $t$ is $-B(t\omega)/t$,
one equation for every wavenumber.
\S\ref{sec:generator} makes that equation an operator on the signal.

\subsection{The characterization theorem}\label{sec:characterization-recap}

\emph{The canonical gauge.} Scale covariance (A8) says that dilating space by $\lambda$
relabels the scales by a bijection $S_\lambda$, and it says nothing about what that bijection
is. The line paper shows (its Proposition~\linenum{prop:canonical-gauge}) that
$c(\lambda) := S_\lambda 1$ is a continuous strictly increasing bijection of $(0,\infty)$ with
$c(1) = 1$, and defines the \emph{canonical gauge} as its inverse, $\chi := c^{-1}$, extended
by $\chi(0) = 0$. In the coordinate $\chi(t)$ the relabelling is dilation itself,
$\chi(S_\lambda t) = \lambda\,\chi(t)$, and the accumulated exponent is
$G(t,\omega) = F(\chi(t)\,\omega)$ with $F := G(1,\cdot)$. A family is \emph{in canonical
gauge} when $\chi$ is the identity, that is when the kernel at scale $\lambda t$ is the
dilation by $\lambda$ of the kernel at scale $t$. Every admissible family is brought to
canonical gauge by reparametrizing its scale, and the gauge is unique once $\chi(1) = 1$ is
fixed. In canonical gauge the family is generated by one law: the kernel at scale $t$ is
the law of $t X_1$, where $X_1$ is the random displacement from scale $0$ to scale $1$.

% shared with blueprint thm:main-characterization
\begin{theorem}[main theorem, Theorem~7.3 of the line paper]\label{thm:main-characterization}
A family $(\Phis{s}{t})_{0 \le s \le t}$ satisfies (A1)--(A8) and (ND) \emph{if and only if}
there are an increasing bijection $\chi$ of $[0,\infty)$, the gauge, and a function $F$ of the
form \ref{eq:sd-profile} with $F \not\equiv 0$, such that
\[
  \Phis{s}{t} f = \mu_{s,t} * f, \qquad
  \hat\mu_{s,t}(\omega) = \exp\bigl[-\bigl(F(\chi(t)\,\omega) - F(\chi(s)\,\omega)\bigr)\bigr].
\]
In the canonical gauge the kernel at scale $t$ is the symmetric probability measure $\phi_t$
with $\hat\phi_t = e^{-F(t\,\cdot)}$, and the pair $(\chi,F)$ is unique up to the normalization
$\chi(1) = 1$.
\end{theorem}

So an admissible family is, in its canonical gauge, an admissible pair $(a,k)$ and nothing
else: a Gaussian coefficient and a nonincreasing displacement profile. The word \emph{admissible},
said of an exponent $F$ in this paper, means of the form \ref{eq:sd-profile} with an
admissible pair, and the \emph{admissible cone} of \S\ref{sec:cone} is the set of such
exponents. The theorem rests on two cited facts, Bochner's theorem in the forward direction
and the symmetric \Levy--Khintchine representation at its converse and uniqueness clauses
(\S\ref{sec:machine-checked}).

The semigroup class is the case where the stages depend only on the difference of their
scales, and the theorem recovers it.

% shared with blueprint cor:semigroup-case
\begin{corollary}[recovering the semigroup case, Corollary~7.4 of the line paper]\label{cor:semigroup-case}
Assume in addition that the family is one-parameter, that is, $\Phis{s}{t}$ depends only on
$t - s$. Then, after the normalization $g_{0,1}(1) = 1$,
\[
  F(\omega) = |\omega|^\alpha \quad \text{for some } 0 < \alpha \le 2 :
\]
the kernels are exactly the symmetric stable densities, which are the members of the extended
semigroup class of \sscite{pauwels1995extended} with positive kernels. Here $\alpha = 2$ is the
Gaussian, and $0 < \alpha < 2$ has Gaussian coefficient $0$ and profile
$k(x) = C_\alpha^{-1}x^{-\alpha}$ with
$C_\alpha := \int_0^\infty (1 - \cos u)\,u^{-1-\alpha}du \in (0,\infty)$.
\end{corollary}

Two consequences of the theorem are used throughout: the admissible exponents form a cone, and
the kernels have densities.

% shared with blueprint lem:admissible-cone
\begin{lemma}[the admissible exponents form a convex cone, Lemma~7.6 of the line paper]\label{lem:admissible-cone}
If $F_1, F_2$ are admissible in the sense of Theorem~\ref{thm:main-characterization}, of the
form \ref{eq:sd-profile} with Gaussian coefficients $a_1, a_2 \ge 0$ and nonincreasing profiles
$k_1, k_2$, then so are $F_1 + F_2$, with data $(a_1+a_2, k_1+k_2)$, and $c\,F_1$ for $c \ge 0$,
with data $(c\,a_1, c\,k_1)$. So the admissible exponents form a convex cone, and the
correspondence with the pairs $(a,k)$ is linear.
\end{lemma}

% shared with blueprint prop:kernel-regularity
\begin{proposition}[regularity of the kernels from the origin, Proposition~7.9 of the line paper]\label{prop:kernel-regularity}
Let $(\Phis{s}{t})$ satisfy (A1)--(A8) and (ND). Then for every $t > 0$ the kernel $\mu_{0,t}$
is absolutely continuous, and its density is unimodal with mode at the origin.
\end{proposition}

The proposition rests on two cited facts about self-decomposable laws: absolute continuity
(\sscite{sato1999levy}, Thm.~27.13, p.~181) and unimodality
(\sscite{yamazato1978unimodality}, Thm.~1, p.~523). The theorem does not use them. It is
stated for measures, so that it rests only on Bochner's theorem and the \Levy--Khintchine
representation. A reader who wants the density reads the theorem and the proposition together.

\subsection{What the line paper takes on trust}\label{sec:line-trust}

The line paper rests its theorems on cited facts about positive definite functions and
self-decomposable laws, and this paper inherits them where it reads its statements.
\S\ref{sec:machine-checked} lists them, with what this paper adds, and the page each rests on
is printed before the statement that uses it.

\subsection{One proposition the line paper did not print}\label{sec:laplace-uniqueness}

The one result of the line paper's preliminaries that is stated here and not cited is the
uniqueness of the Laplace transform of a measure on the half-line, in the form that allows the
measure to be infinite. It is used twice. In \S\ref{sec:corners} it identifies the mixing
measure of a profile that is a mixture of exponentials, its Thorin measure. In
\S\ref{sec:bridge} it rules out a causal ancestor. The result is elementary once finiteness of
the transform is used to reduce to finite measures, and its proof uses finiteness at \emph{one}
point of $(\tau_0,\infty)$, which is weaker than the hypothesis stated.

% shared with blueprint prop:laplace-uniqueness-locally-finite
\begin{proposition}[uniqueness of the Laplace transform of a locally finite measure]
\label{prop:laplace-uniqueness-locally-finite}
Let $m, m'$ be Borel measures on $(0,\infty)$, not necessarily finite, and suppose that
$\int e^{-\tau u}\,m(du) = \int e^{-\tau u}\,m'(du) < \infty$ for every $\tau$ in some
interval $(\tau_0,\infty)$. Then $m = m'$.
\end{proposition}

\begin{proof}
Fix $\tau_1 > \tau_0$. Both transforms being finite at $\tau_1$, the measures
$e^{-\tau_1 u}m(du)$ and $e^{-\tau_1 u}m'(du)$ are finite, and they have equal Laplace
transforms on $(0,\infty)$ after the shift $\tau \mapsto \tau + \tau_1$; so it is enough to
treat finite measures. For finite $m$, the map $u \mapsto e^{-u}$ carries $(0,\infty)$
homeomorphically onto $(0,1)$, and $\int e^{-n u}\,m(du)$, $n \in \NN$, is the $n$-th moment
of the image measure on $[0,1]$. A finite measure on a bounded interval is determined by its
moments, the polynomials being dense in the continuous functions there. So the image measures
agree, and so do $m$ and $m'$.
\end{proof}
